\documentclass{article}
\usepackage[T1]{fontenc}
\usepackage{lmodern}
\usepackage{graphicx}
\usepackage{amsmath,amssymb}
\usepackage[a4paper,margin=2cm]{geometry}
\usepackage[absolute,overlay]{textpos}
\setlength{\TPHorizModule}{1mm}
\setlength{\TPVertModule}{1mm}
\usepackage[indonesian]{babel}
\usepackage{hyperref}
\setlength{\emergencystretch}{2em}
% unit: Halaman 9

\begin{document}

% === Halaman 9 (llm) ===

\begin{itemize}
    \item Kurva eliptik $y^2 = x^3 + ax + b$ terdefinisi untuk $x, y \in \mathbb{R}$
    \item Didefinisikan sebuah titik bernama titik $O(x, \infty)$, yaitu titik pada \textit{infinity}.
    \item Titik-titik $P(x, y)$ pada kurva eliptik bersama operasi $+$ membentuk sebuah grup.
    \begin{itemize}
        \item Himpunan $G$ : semua titik $P(x, y)$ pada kurva eliptik
        \item Operasi biner : $+$
    \end{itemize}
    \item Penjelasan kenapa kurva eliptik membentuk sebuah grup dijelaskan pada \textit{slide-slide} berikut ini.
\end{itemize}

\end{document}
